% ----------------------------------------------------------
% Appendix Main Heading (ONLY this goes to TOC)
% ----------------------------------------------------------

\clearpage
\chapter*{Appendix}
\addcontentsline{toc}{chapter}{Appendix}

\vspace{1cm}

% ----------------------------------------------------------
% Appendix I - Code Snippets
% ----------------------------------------------------------

\section*{Appendix I -- Code Snippets}

\vspace{0.5cm}

The following listings present representative source-code excerpts from the MetaGPT implementation,
covering backend API configuration, agent orchestration, storage, frontend interaction, and
retrieval-augmented generation.

\vspace{0.5cm}

\begin{center}
\begin{minipage}{\linewidth}
\begin{lstlisting}[language=Python, caption={FastAPI application configuration and middleware setup}, label={lst:fastapi_setup}]
def create_app() -> FastAPI:
    settings = get_settings()
    app = FastAPI(
        title=settings.app_name,
        version=settings.app_version,
        lifespan=lifespan,
    )
    # Configure CORS
    app.add_middleware(
        CORSMiddleware,
        allow_origins=settings.cors_origins,
        allow_credentials=True,
        allow_methods=["*"],
        allow_headers=["*"],
    )
    # Include API routes
    app.include_router(api_router, prefix=settings.api_prefix)
    return app
\end{lstlisting}
\end{minipage}
\end{center}

\vspace{0.5cm}

\begin{center}
\begin{minipage}{\linewidth}
\begin{lstlisting}[language=Python, caption={FastAPI endpoint for project creation}, label={lst:backend_create_project}]
@router.post("", response_model=Project, status_code=status.HTTP_201_CREATED)
async def create_project(
    request: ProjectCreate,
    user: User = Depends(get_current_user),
) -> Project:
    if not user.has_credits():
        raise HTTPException(
            status_code=status.HTTP_403_FORBIDDEN,
            detail="You've used all your free credits. Upgrade to premium.",
        )
    service = PipelineService()
    project = await service.create_project(request.prompt, user_id=str(user.id))
    user.credits_used += 1
    await user.save()
    return project
\end{lstlisting}
\end{minipage}
\end{center}

\vspace{0.5cm}

\begin{center}
\begin{minipage}{\linewidth}
\begin{lstlisting}[language=Python, caption={LangGraph orchestrator pipeline configuration}, label={lst:orchestrator_graph}]
def _build_graph(self) -> StateGraph:
    workflow = StateGraph(PipelineState)
    # Add nodes for each agent
    workflow.add_node("manager", self._run_manager)
    workflow.add_node("architect", self._run_architect)
    workflow.add_node("engineer", self._run_engineer)
    workflow.add_node("qa", self._run_qa)

    # Define deterministic linear flow
    workflow.set_entry_point("manager")
    workflow.add_edge("manager", "architect")
    workflow.add_edge("architect", "engineer")
    workflow.add_edge("engineer", "qa")
    workflow.add_edge("qa", END)
    
    return workflow.compile()
\end{lstlisting}
\end{minipage}
\end{center}

\vspace{0.5cm}

\begin{center}
\begin{minipage}{\linewidth}
\begin{lstlisting}[language=Python, caption={Engineer Agent building its prompt with context from earlier pipeline stages}, label={lst:engineer_prompt}]
def _build_prompt(self, **inputs: Any) -> str:
    architect_output: ArchitectOutput = inputs.get("architect_output")
    manager_output: ManagerOutput = inputs.get("manager_output")

    # Format SOP sections
    constraints_text = self._format_sop_section(self.sop.constraints)
    quality_text = self._format_sop_section(self.sop.quality_checklist)

    # Format architecture details
    file_structure = self._format_file_structure(architect_output)
    components = self._format_components(architect_output)
    requirements = self._format_requirements(manager_output)

    return f"""# Engineer Agent - Code Generation
## Your Role
{self.sop.role}

## Constraints
{constraints_text}

## Project: {manager_output.project_name}
{manager_output.project_description}

## File Structure:
{file_structure}

## Your Task
Implement ALL files listed in the file structure with production-ready code.
"""
\end{lstlisting}
\end{minipage}
\end{center}

\vspace{0.5cm}

\begin{center}
\begin{minipage}{\linewidth}
\begin{lstlisting}[language=Python, caption={Creating and persisting a project document using Beanie ODM}, label={lst:mongo_create}]
async def create(self, project_id: str, prompt: str, user_id: str = "") -> Project:
    doc = ProjectDocument(
        project_id=project_id,
        user_id=user_id,
        prompt=prompt,
        state=ProjectState(
            pipeline_status=PipelineStatus(
                stage=PipelineStage.PENDING,
                progress=0,
            )
        ),
        created_at=datetime.utcnow(),
        updated_at=datetime.utcnow(),
    )
    await doc.insert()
    return self._to_project(doc)
\end{lstlisting}
\end{minipage}
\end{center}

\vspace{0.5cm}

\begin{center}
\begin{minipage}{\linewidth}
\begin{lstlisting}[language=JavaScript, caption={Client-side form submission handler in Next.js}, label={lst:frontend_submit}]
const handleSubmit = async (e: React.FormEvent) => {
    e.preventDefault();
    if (!prompt.trim() || isLoading) return;
    if (!token) {
        router.push("/signin");
        return;
    }
    setIsLoading(true);
    try {
        const { createProject } = await import("@/lib/api");
        const project = await createProject(prompt.trim());
        router.push(`/project/${project.id}`);
    } catch (error: unknown) {
        setSubmitError("Failed to create project");
        setIsLoading(false);
    }
};
\end{lstlisting}
\end{minipage}
\end{center}

\vspace{0.5cm}

\begin{center}
\begin{minipage}{\linewidth}
\begin{lstlisting}[language=JavaScript, caption={Dynamically building a file tree structure from a flat list of file paths}, label={lst:frontend_file_tree}]
function buildTreeFromFiles(files: { file_path: string }[]): FileTreeNode | null {
  if (files.length === 0) return null;

  const root: FileTreeNode = { name: "files", path: "/", type: "directory", children: [], };

  for (const file of files) {
    const parts = file.file_path.split("/").filter(Boolean);
    let current = root;

    for (let i = 0; i < parts.length; i++) {
      const part = parts[i];
      const isFile = i === parts.length - 1;
      const path = parts.slice(0, i + 1).join("/");
      let child = current.children?.find((c) => c.name === part);

      if (!child) {
        child = { name: part, path, type: isFile ? "file" : "directory", children: isFile ? [] : [], };
        current.children = current.children || [];
        current.children.push(child);
      }
      current = child;
    }
  }
  return root;
}
\end{lstlisting}
\end{minipage}
\end{center}

\vspace{0.5cm}

\begin{center}
\begin{minipage}{\linewidth}
\begin{lstlisting}[language=JavaScript, caption={Chat handler supporting Server-Sent Events (SSE) and background indexing}, label={lst:chat_submit}]
const handleSubmit = async (e: React.FormEvent) => {
    e.preventDefault();
    if (!message.trim() || isLoading) return;

    const userMessage = message.trim();
    setMessage("");
    setIsLoading(true);

    try {
        const response = await sendChatMessage(
            projectId,
            userMessage,
            autoMode ? null : selectedModel,
        );

        addChatMessage({
            ...response.message,
            files_referenced: response.files_referenced,
            indexing_status: response.indexing_status,
        });

        // If indexing is pending, trigger it live
        if (response.indexing_status === "pending") {
            indexProjectFiles(projectId, response.files_modified);
        }
    } finally {
        setIsLoading(false);
    }
};
\end{lstlisting}
\end{minipage}
\end{center}

\vspace{0.5cm}

\begin{center}
\begin{minipage}{\linewidth}
\begin{lstlisting}[language=Python, caption={Configuring Gemini for structured JSON output via LangChain}, label={lst:gemini_structured}]
def get_llm_with_structured_output(
    output_schema: type[T],
    temperature: float | None = None,
    max_tokens: int | None = None,
    model: str | None = None,
) -> BaseChatModel:
    base_llm = get_llm(temperature=temperature, max_tokens=max_tokens, model=model)
    return base_llm.with_structured_output(output_schema)
\end{lstlisting}
\end{minipage}
\end{center}

\vspace{0.5cm}

\begin{center}
\begin{minipage}{\linewidth}
\begin{lstlisting}[language=Python, caption={Retrieving relevant codebase chunks from the vector store}, label={lst:retriever_search}]
async def retrieve(self, project_id: str, query: str, k: int | None = None) -> list[RetrievedChunk]:
    if not self.settings.rag_enabled:
        return []

    k = k or self.settings.rag_top_k
    try:
        vector_store = get_vector_store(project_id)
        # Similarity search with scores
        results = vector_store.similarity_search_with_relevance_scores(query, k=k)

        chunks = []
        for doc, score in results:
            chunks.append(
                RetrievedChunk(
                    content=doc.page_content,
                    file_path=doc.metadata.get("file_path", "unknown"),
                    language=doc.metadata.get("language", "text"),
                    relevance_score=score,
                )
            )
        return chunks
    except Exception as e:
        logger.error(f"Retrieval error for project {project_id}: {e}")
        return []
\end{lstlisting}
\end{minipage}
\end{center}

\vspace{0.5cm}

\begin{center}
\begin{minipage}{\linewidth}
\begin{lstlisting}[language=Python, caption={BaseAgent: core LLM invocation loop with structured output and execution statistics}, label={lst:base_agent_execute}]
async def execute(self, **inputs: Any) -> T:
    start_time = time.time()

    # Build the prompt from SOP template
    prompt = self._build_prompt(**inputs)

    # Get LLM configured for structured output
    llm = get_llm_with_structured_output(self.output_schema)

    # Create messages
    messages = [
        SystemMessage(content=self.sop.role),
        HumanMessage(content=prompt),
    ]

    # Execute and get structured response
    try:
        result = await llm.ainvoke(messages)

        self._execution_stats = {
            "execution_time_ms": int((time.time() - start_time) * 1000),
            "status": "success",
            "agent_name": self.name,
        }
        return result

    except Exception as e:
        self._execution_stats = {
            "execution_time_ms": int((time.time() - start_time) * 1000),
            "status": "error",
            "error": str(e),
            "agent_name": self.name,
        }
        raise
\end{lstlisting}
\end{minipage}
\end{center}

\vspace{0.5cm}

\begin{center}
\begin{minipage}{\linewidth}
\begin{lstlisting}[language=Python, caption={LangGraph PipelineState TypedDict and initial state factory}, label={lst:pipeline_state}]
class PipelineState(TypedDict, total=False):
    """State that flows through the LangGraph pipeline."""

    # Input
    project_id: str
    user_prompt: str
    context: str

    # Agent outputs
    manager_output:  ManagerOutput  | None
    architect_output: ArchitectOutput | None
    engineer_output:  EngineerOutput  | None
    qa_output:        QAOutput        | None

    # Pipeline tracking
    current_stage: Literal[
        "pending", "manager", "architect", "engineer", "qa", "completed", "error"
    ]
    progress: float
    error: str | None
    execution_log: list[dict[str, Any]]


def create_initial_state(project_id: str, user_prompt: str, context: str = "") -> PipelineState:
    """Create initial pipeline state from user input."""
    return PipelineState(
        project_id=project_id,
        user_prompt=user_prompt,
        context=context,
        manager_output=None,
        architect_output=None,
        engineer_output=None,
        qa_output=None,
        current_stage="pending",
        progress=0.0,
        error=None,
        started_at=datetime.utcnow().isoformat(),
        completed_at=None,
        execution_log=[],
    )
\end{lstlisting}
\end{minipage}
\end{center}

\vspace{0.5cm}

\begin{center}
\begin{minipage}{\linewidth}
\begin{lstlisting}[language=Python, caption={ManagerAgent: SOP-driven prompt construction for requirements extraction}, label={lst:manager_prompt}]
class ManagerAgent(BaseAgent[ManagerOutput]):
    """Manager Agent - Converts user prompts to structured requirements."""

    def __init__(self):
        super().__init__(sop=MANAGER_SOP, output_schema=ManagerOutput)

    @property
    def name(self) -> str:
        return "ManagerAgent"

    def _build_prompt(self, **inputs: Any) -> str:
        user_prompt = inputs.get("user_prompt", "")
        context = inputs.get("context", "No additional context provided.")

        constraints_text = self._format_sop_section(self.sop.constraints)
        quality_text = self._format_sop_section(self.sop.quality_checklist)

        return self.sop.prompt_template.format(
            role=self.sop.role,
            objective=self.sop.objective,
            constraints=constraints_text,
            quality_checklist=quality_text,
            user_prompt=user_prompt,
            context=context,
        )


async def run_manager_agent(user_prompt: str, context: str = "") -> ManagerOutput:
    """Convenience function to run the Manager Agent."""
    agent = ManagerAgent()
    return await agent.execute(user_prompt=user_prompt, context=context)
\end{lstlisting}
\end{minipage}
\end{center}

\vspace{0.5cm}

\begin{center}
\begin{minipage}{\linewidth}
\begin{lstlisting}[language=Python, caption={JWT authentication: token creation and FastAPI dependency for current user}, label={lst:jwt_auth}]
def create_access_token(user_id: str, email: str) -> str:
    """Create a JWT access token for a user."""
    settings = get_settings()
    expire = datetime.utcnow() + timedelta(hours=settings.jwt_expiration_hours)
    payload = {"sub": user_id, "email": email, "exp": expire}
    return jwt.encode(payload, settings.jwt_secret, algorithm=settings.jwt_algorithm)


async def get_current_user(
    credentials: Annotated[HTTPAuthorizationCredentials, Depends(security)],
) -> User:
    """FastAPI dependency: validates JWT and returns the User from MongoDB."""
    settings = get_settings()
    token = credentials.credentials

    credentials_exception = HTTPException(
        status_code=status.HTTP_401_UNAUTHORIZED,
        detail="Invalid or expired token",
        headers={"WWW-Authenticate": "Bearer"},
    )

    try:
        payload = jwt.decode(token, settings.jwt_secret, algorithms=[settings.jwt_algorithm])
        user_id: str | None = payload.get("sub")
        if user_id is None:
            raise credentials_exception
    except JWTError:
        raise credentials_exception

    try:
        user = await User.get(PydanticObjectId(user_id))
    except Exception:
        raise credentials_exception

    if user is None:
        raise credentials_exception

    return user
\end{lstlisting}
\end{minipage}
\end{center}

\vspace{0.5cm}

\begin{center}
\begin{minipage}{\linewidth}
\begin{lstlisting}[language=Python, caption={CodebaseIndexer: language-aware file chunking and vector store ingestion}, label={lst:indexer_project}]
def _get_splitter(self, file_path: str) -> RecursiveCharacterTextSplitter:
    """Get a language-aware text splitter for the file type."""
    ext = Path(file_path).suffix.lower()
    language = _LANGUAGE_MAP.get(ext)

    if language:
        return RecursiveCharacterTextSplitter.from_language(
            language=language,
            chunk_size=self.settings.rag_chunk_size,
            chunk_overlap=self.settings.rag_chunk_overlap,
        )
    return RecursiveCharacterTextSplitter(
        chunk_size=self.settings.rag_chunk_size,
        chunk_overlap=self.settings.rag_chunk_overlap,
        separators=["\n\n", "\n", " ", ""],
    )


async def index_project(self, project_id: str) -> dict:
    """Index all files in a project into the vector store."""
    file_paths = await self.file_store.list_files(project_id)
    documents = []

    for file_path in file_paths:
        if self._should_skip(file_path):
            continue
        file_data = await self.file_store.read_file(project_id, file_path)
        if not file_data or not file_data.content.strip():
            continue

        splitter = self._get_splitter(file_path)
        for i, chunk in enumerate(splitter.split_text(file_data.content)):
            documents.append(Document(
                page_content=chunk,
                metadata={
                    "file_path": file_path,
                    "language": file_data.language,
                    "chunk_index": i,
                    "project_id": project_id,
                    "indexed_at": datetime.utcnow().isoformat(),
                },
            ))

    vector_store = get_vector_store(project_id)
    vector_store.add_documents(documents)
    return {"files_indexed": len(set(d.metadata["file_path"] for d in documents)),
            "chunks_created": len(documents)}
\end{lstlisting}
\end{minipage}
\end{center}

\vspace{0.5cm}

\begin{center}
\begin{minipage}{\linewidth}
\begin{lstlisting}[language=Python, caption={ChatService: RAG-driven file selection and single-LLM-call edit flow}, label={lst:chat_service}]
async def process_message(self, project_id: str, request: ChatRequest) -> ChatResponse:
    """
    Process a chat message to produce targeted file edits.
    Flow: RAG retrieval -> read file contents -> single LLM call -> write changes.
    """
    all_file_paths = await self.file_store.list_files(project_id)

    # Build file context: all files for small projects, RAG-selected for large ones
    if len(all_file_paths) <= 30:
        file_context = await self._read_all_files(project_id, all_file_paths)
    elif self.settings.rag_enabled:
        relevant = await self.retriever.retrieve_files(project_id, request.message, k=15)
        file_context = "\n\n".join(
            f"### {f['file_path']} ({f['language']})\n```{f['language']}\n{f['content']}\n```"
            for f in relevant
        )
    else:
        file_context = await self._read_all_files(project_id, all_file_paths)

    # Single LLM call for targeted edits
    llm_result = await self._call_llm(request.message, file_context, all_file_paths)

    # Write modified files to disk and update MongoDB
    files_modified = []
    for file_data in llm_result.get("files", []):
        spec = GeneratedFileSpec(**file_data)
        await self.file_store.write_file(project_id, spec)
        files_modified.append(spec.file_path)

    if files_modified:
        await self._update_project_state(project_id, files_modified)

    return ChatResponse(
        message=ChatMessage(role="assistant",
                            content=llm_result.get("summary", "Done.")),
        files_modified=files_modified,
        indexing_status="pending" if files_modified and self.settings.rag_enabled else None,
    )
\end{lstlisting}
\end{minipage}
\end{center}

\vspace{0.5cm}

\begin{center}
\begin{minipage}{\linewidth}
\begin{lstlisting}[language=Python, caption={MongoDB initialisation with Motor and Beanie ODM}, label={lst:db_init}]
async def init_db() -> None:
    """Initialize the MongoDB connection and Beanie ODM."""
    from app.models.user import User
    from app.models.project import ProjectDocument

    settings = get_settings()
    client = AsyncIOMotorClient(settings.mongodb_uri)
    db = client[settings.mongodb_db]

    await init_beanie(
        database=db,
        document_models=[User, ProjectDocument],
    )
\end{lstlisting}
\end{minipage}
\end{center}
